\begingroup
\renewcommand{\arraystretch}{1.45}
\fancyhead[L]{\small DETERMINATION OF ATTRACTION}
\Needspace{32\baselineskip}
\AttractionArticle{3}
The task now comes down to reducing the radical expression to a simpler form by introducing another variable in place of $E$. For this purpose we shall set
\[
\begin{aligned}
&\cos E=\frac{\alpha+\alpha'\cos T+\alpha''\sin T}{\gamma+\gamma'\cos T+\gamma''\sin T},\\[6pt]
&
\sin E=\frac{\beta+\beta'\cos T+\beta''\sin T}{\gamma+\gamma'\cos T+\gamma''\sin T}
\end{aligned}
\]
Here, however, the nine coefficients $\alpha,\alpha',\alpha''$ etc. are obviously not entirely arbitrary, but must satisfy certain conditions, which must be investigated first of all. We first observe that the substitution remains the same if all the coefficients are multiplied by the same factor, so that, without loss of generality, one of them could be assigned a fixed value, e. g. one could set $\gamma=\nobreak{}1$: nevertheless, for the sake of elegance, let all nine remain indeterminate for a while. Further, we point out that values must be ex\SourceLeaf{344}%
cluded for which $\alpha,\alpha',\alpha''$ or $\beta,\beta',\beta''$ would be respectively proportional to $\gamma,\gamma'.\gamma''$: otherwise $E$ would no longer remain indeterminate. Therefore $\gamma'\alpha''-\gamma''\alpha',\gamma''\alpha-\gamma\alpha'',\gamma\alpha'-\gamma'\alpha$ cannot all vanish at once.

\Needspace{20\baselineskip}
The coefficients $\alpha,\alpha',\alpha''$ etc. must obviously be such that the following holds identically:
\[
\left.\begin{aligned}
 &(\alpha+\alpha'\cos T+\alpha''\sin T)^2\\[5pt]
+&(\beta+\beta'\cos T+\beta''\sin T)^2\\[5pt]
-&(\gamma+\gamma'\cos T+\gamma''\sin T)^2
\end{aligned}\right\}=0
\]
whence this function must necessarily have the form
\[
\begin{aligned}
&k(\cos T^2+\sin T^2-1)
\end{aligned}
\]
From this we obtain six equations of condition
\[
\left.\begin{aligned}
-\alpha\alpha-\beta\beta+\gamma\gamma&=k\\[5pt]
-\alpha'\alpha'-\beta'\beta'+\gamma'\gamma'&=-k\\[5pt]
-\alpha''\alpha''-\beta''\beta''+\gamma''\gamma''&=-k\\[5pt]
-\alpha'\alpha''-\beta'\beta''+\gamma'\gamma''&=0\\[5pt]
-\alpha''\alpha-\beta''\beta+\gamma''\gamma&=0\\[5pt]
-\alpha\alpha'-\beta\beta'+\gamma\gamma'&=0
\end{aligned}\right\}\quad\mathrm{(I)}
\]
\Needspace{27\baselineskip}
Several other equations depend on these, and it will be worthwhile to develop them. On setting, for brevity,
\[
\begin{aligned}
&\alpha\beta'\gamma''+\alpha'\beta''\gamma+\alpha''\beta\gamma'
-\alpha\beta''\gamma'-\alpha'\beta\gamma''-\alpha''\beta'\gamma
\\[6pt]
&\qquad{}=\varepsilon\quad\mathinner{.\,.\,.\,.\,.\,.\,.\,.}\quad\text{(II)}
\end{aligned}
\]
the following nine are easily derived by combining equations (I):
\[
\left.\begin{aligned}
\varepsilon\alpha&=-k(\beta'\gamma''-\gamma'\beta'')\\[5pt]
\varepsilon\beta&=-k(\gamma'\alpha''-\alpha'\gamma'')\\[5pt]
\varepsilon\gamma&=+k(\alpha'\beta''-\beta'\alpha'')\\[5pt]
\varepsilon\alpha'&=+k(\beta''\gamma-\gamma''\beta)\\[5pt]
\varepsilon\beta'&=+k(\gamma''\alpha-\alpha''\gamma)\\[5pt]
\varepsilon\gamma'&=-k(\alpha''\beta-\beta''\alpha)\\[5pt]
\varepsilon\alpha''&=+k(\beta\gamma'-\gamma\beta')\\[5pt]
\varepsilon\beta''&=+k(\gamma\alpha'-\alpha\gamma')\\[5pt]
\varepsilon\gamma''&=-k(\alpha\beta'-\beta\alpha')
\end{aligned}\right\}\quad\mathrm{(III)}
\]
\Needspace{38\baselineskip}
\SourceLeaf{345}
From the first three of these equations we again deduce the following:
\[
\begin{aligned}
&\varepsilon\alpha(\beta'\gamma''-\gamma'\beta'')+\varepsilon\beta(\gamma'\alpha''-\alpha'\gamma'')+\varepsilon\gamma(\alpha'\beta''-\beta'\alpha'')\\[6pt]
&=-k(\beta'\gamma''-\gamma'\beta'')^2-k(\gamma'\alpha''-\alpha'\gamma'')^2\\[6pt]
&\qquad{}+k(\alpha'\beta''-\beta'\alpha'')^2
\end{aligned}
\]
to which the following is equivalent:
\[
\begin{aligned}
&\varepsilon\varepsilon=k(-\alpha'\alpha'-\beta'\beta'+\gamma'\gamma')(-\alpha''\alpha''-\beta''\beta''+\gamma''\gamma'')\\[6pt]
&\qquad{}-k(-\alpha'\alpha''-\beta'\beta''+\gamma'\gamma'')^2
\end{aligned}
\]
which, with the aid of equations 2, 3, 4 in (I), becomes the following:
\[
\begin{aligned}
&\varepsilon\varepsilon=k^3\;\ldotp\;\ldotp\;\ldotp\;\ldotp\;\ldotp\;\ldotp\;\ldotp\;\ldotp\;\ldotp\;\ldotp\;\ldotp\;\text{(IV)}
\end{aligned}
\]
The following are derived just as easily from equations (I):
\[
\left.\begin{aligned}
(\beta'\gamma''-\gamma'\beta'')^2&=-k(k-\alpha'\alpha'-\alpha''\alpha'')\\[5pt]
(\gamma'\alpha''-\alpha'\gamma'')^2&=-k(k-\beta'\beta'-\beta''\beta'')\\[5pt]
(\alpha'\beta''-\beta'\alpha'')^2&=+k(k+\gamma'\gamma'+\gamma''\gamma'')\\[5pt]
(\beta''\gamma-\gamma''\beta)^2&=+k(k+\alpha\alpha-\alpha''\alpha'')\\[5pt]
(\gamma''\alpha-\alpha''\gamma)^2&=+k(k+\beta\beta-\beta''\beta'')\\[5pt]
(\alpha''\beta-\beta''\alpha)^2&=-k(k-\gamma\gamma+\gamma''\gamma'')\\[5pt]
(\beta\gamma'-\gamma\beta')^2&=+k(k+\alpha\alpha-\alpha'\alpha')\\[5pt]
(\gamma\alpha'-\alpha\gamma')^2&=+k(k+\beta\beta-\beta'\beta')\\[5pt]
(\alpha\beta'-\beta\alpha')^2&=-k(k-\gamma\gamma+\gamma'\gamma')
\end{aligned}\right\}\quad\text{(V)}
\]

\Needspace{14\baselineskip}
As an example we write out the expansion of the first, on whose model the others will readily be formed. Namely, equations 4, 2, 3 in (I) give
\[
\begin{aligned}
&(\gamma'\gamma''-\beta'\beta'')^2-(\gamma'\gamma'-\beta'\beta')(\gamma''\gamma''-\beta''\beta'')\\[6pt]
&\qquad{}=\alpha'\alpha'\alpha''\alpha''-(\alpha'\alpha'-k)(\alpha''\alpha''-k)
\end{aligned}
\]
which equation, when expanded, immediately gives the first in (V) itself.

\Needspace{26\baselineskip}
From these equations (V) we conclude that the value $k=0$ is not admissible in our investigation; for then all nine quantities $\beta'\gamma''-\gamma'\beta''$ etc. would necessarily vanish, i. e. the coefficients $\alpha$, $\alpha'$, $\alpha''$ would become proportional both to $\beta$, $\beta'$, $\beta''$ and to $\gamma$, $\gamma'$, $\gamma''$. Hence also, by equation IV, the quantity $\varepsilon$ cannot vanish; consequently $k$ must necessarily be a positive quantity, provided that all the coefficients $\alpha$, $\alpha'$, $\alpha''$ etc. are to be real. Combining the first three equations in (III) with the first three in (V) gives the following new equations, which obviously depend on $k$ being nonzero:
\SourceLeaf{346}
\[
\left.\begin{aligned}
\alpha\alpha-\alpha'\alpha'-\alpha''\alpha''&=-k\\[5pt]
\beta\beta-\beta'\beta'-\beta''\beta''&=-k\\[5pt]
\gamma\gamma-\gamma'\gamma'-\gamma''\gamma''&=+k
\end{aligned}\right\}\quad\text{(VI)}
\]
\Needspace{20\baselineskip}
Combining the remaining ones would produce the same equations. Finally, to these we add the following three:
\[
\left.\begin{aligned}
\beta\gamma-\beta'\gamma'-\beta''\gamma''&=0\\[5pt]
\gamma\alpha-\gamma'\alpha'-\gamma''\alpha''&=0\\[5pt]
\alpha\beta-\alpha'\beta'-\alpha''\beta''&=0
\end{aligned}\right\}\quad\text{(VII)}
\]
which are easily derived from equations III; e. g. the second, fifth and eighth give:
\[
\begin{aligned}
&\varepsilon\beta\gamma-\varepsilon\beta'\gamma'-\varepsilon\beta''\gamma''=-k\gamma(\gamma'\alpha''-\alpha'\gamma'')\\[6pt]
&\qquad{}-k\gamma'(\gamma''\alpha-\alpha''\gamma)-k\gamma''(\gamma\alpha'-\alpha\gamma')=0
\end{aligned}
\]
These equations too obviously depend on the exclusion of the value $k=0$\HistoricalNote{25}.

\Needspace{36\baselineskip}
Since, as we have already pointed out above, all the coefficients $\alpha$, $\alpha'$, $\alpha''$ etc. may be multiplied by the same factor, whereby the value of $k$ will be multiplied by the square of that factor, from now on we shall always suppose
\[
\begin{aligned}
&k=1
\end{aligned}
\]
in which case there will necessarily also be either $\varepsilon=+1$ or $\varepsilon=-1$. It is therefore evident that the nine coefficients $\alpha$, $\alpha'$, $\alpha''$ etc., between which there are six equations of condition, must be reducible to three quantities independent of one another, which is accomplished most conveniently by means of three angles as follows:
\[
\begin{aligned}
&\alpha=\cos L\operatorname{tang}N\\[6pt]
&\beta=\sin L\operatorname{tang}N\\[6pt]
&\gamma=\sec N\\[6pt]
&\alpha'=\cos L\cos M\sec N\pm\sin L\sin M\\[6pt]
&\beta'=\sin L\cos M\sec N\mp\cos L\sin M\\[6pt]
&\gamma'=\cos M\operatorname{tang}N\\[6pt]
&\alpha''=\cos L\sin M\sec N\mp\sin L\cos M\\[6pt]
&\beta''=\sin L\sin M\sec N\pm\cos L\cos M\\[6pt]
&\gamma''=\sin M\operatorname{tang}N
\end{aligned}
\]

\SourceLeaf{347}\noindent where the upper of the ambiguous signs refer to the case $\varepsilon=+1$, and the lower to the case $\varepsilon=-1$. Yet the analytical treatment is for the most part carried out more elegantly without using these angles. Moreover, it would not be difficult to assign a geometrical meaning both to these angles and to the other auxiliary quantities occurring in this investigation; but we leave this interpretation, which is not necessary for our purpose, to the experienced reader to explain.


\clearpage
\endgroup
